% Generated draft: requires mathematical and source review.
\documentclass[11pt]{article}
\usepackage[margin=25mm]{geometry}
\usepackage{fontspec}
\setmainfont{Latin Modern Roman}
\usepackage{amsmath,amssymb,mathtools,unicode-math}
\setmathfont{Latin Modern Math}
\usepackage{enumitem,needspace,xurl,hyperref}
\setlength{\parindent}{0pt}
\setlength{\parskip}{5pt}
\setlength{\emergencystretch}{3em}
\newcommand{\N}{\mathbb{N}}
\newcommand{\Z}{\mathbb{Z}}
\newcommand{\Q}{\mathbb{Q}}
\newcommand{\R}{\mathbb{R}}
\newcommand{\C}{\mathbb{C}}
\newcommand{\F}{\mathbb{F}}
\newcommand{\A}{\mathbb{A}}
\newcommand{\PP}{\mathbb{P}}
\newcommand{\E}{\mathbb{E}}
\newcommand{\eps}{\varepsilon}
\newcommand{\dd}{\,\mathrm{d}}
\begin{document}
\section*{UnsolvedMath ID: 700536}
\textbf{Source problem code:} catalogue entry
\subsection{Definitions and mathematical statement}
For a group \(G\) and \(x\in G\), let \(\operatorname{ord}_G(x)\) denote the order of \(x\), with value \(\infty\) when \(x\) has infinite order. Adopt the convention that every positive integer divides \(\infty\). The group \(G\) is \emph{potent} if, for every \(x\in G\) and every positive integer \(n\) satisfying \(n\mid\operatorname{ord}_G(x)\), there exist a finite group \(Q\) and a surjective homomorphism \(\varphi:G\to Q\) such that
\[
\operatorname{ord}_Q(\varphi(x))=n.
\]
The group \(Q\) and the homomorphism \(\varphi\) may depend on both \(x\) and \(n\).

For groups \(A\) and \(B\), write \(A*B\) for their free product, the group equipped with canonical homomorphisms from \(A\) and \(B\) such that every pair of homomorphisms from these groups to a common group extends uniquely to a homomorphism from \(A*B\).

Is it true that, for every pair of potent groups \(A\) and \(B\), their free product \(A*B\) is potent? No finiteness or finite generation assumption on \(A\) or \(B\) is imposed in the supplied statement.
\subsection{Short English statement}
A group is potent if each element can have any prescribed positive divisor of its order realized as its exact order in a finite quotient; for an element of infinite order, every positive integer is allowed. Must the free product of any two potent groups also be potent?
\subsection{Sources}
\begin{itemize}
\item Supplied local UnsolvedMath registry record, permanent ID \(700536\), exact source problem code \texttt{catalogue entry}; upstream revision \texttt{d371ee917ac9647fbb06cd6fc529deec6f3dda82}. The supplied registry labels the problem \emph{open}; this label has not been independently verified.
\item \emph{The Kourovka Notebook (Unsolved Problems in Group Theory)}, original problem 9.1, 9th issue (1984), proposed by R. B. J. T. Allenby. The supplied background identifies the 21st issue edition (2026), arXiv:1401.0300v48. The collection description names E. I. Khukhro and V. D. Mazurov as editors. These bibliographic details are supplied attribution, not independently verified source consultation.
\end{itemize}
\end{document}
